\chapter{Arquitetura do Sistema}
\label{chap:arquitetura}

O GenJazz é uma \ac{SPA} construída com \textbf{React~18} e \textbf{Vite}, que apresenta a interface dentro de uma moldura de smartphone para deixar claro que a experiência foi pensada para mobile. A arquitetura segue um modelo cliente-servidor simples: toda a lógica de apresentação corre no browser do utilizador, e os dados musicais são fornecidos por uma API REST externa alojada em \texttt{genjazz-api.fly.dev}.

\section{Stack e Organização do Código}
\label{chap2:sec:stack}

A escolha do React deveu-se à sua abordagem declarativa e ao sistema de componentes, que tornam mais fácil gerir os vários estados possíveis de uma interface com tantos elementos dinâmicos. O Vite foi preferido ao Create React App simplesmente por ser mais rápido --- tempos de arranque e hot-reload visivelmente melhores durante o desenvolvimento.

O código está organizado em \texttt{src/screens/} para os ecrãs completos, \texttt{src/components/} para componentes reutilizáveis (como o \texttt{CircleOfFifths}, o \texttt{AudioPlayer} e o \texttt{MeasureGrid}), e \texttt{src/hooks/} para lógica partilhada.

\section{Gestão de Estado e Autenticação}
\label{chap2:sec:estado}

O estado global é gerido por dois contextos React. O \texttt{AuthContext} trata da autenticação: registo, login e dados do utilizador. As palavras-passe são transformadas em hash SHA-256 via \texttt{crypto.subtle} antes de serem guardadas no \texttt{localStorage}. A sessão fica no \texttt{sessionStorage}, pelo que termina automaticamente quando se fecha o separador.

O \texttt{ThemeContext} controla o modo de acessibilidade cromática ativo. Quando o utilizador muda o modo, é adicionado um atributo \texttt{data-colorblind} ao elemento \texttt{<html>}, o que faz com que as variáveis CSS de acento, erro e sucesso sejam substituídas pela paleta correspondente, sem tocar em nenhum componente individualmente.

Toda a lógica da API foi centralizada no hook \texttt{useGenJazz}, que trata do carregamento do catálogo, das seleções do utilizador, da geração de progressões, da conversão para MP3 e da biblioteca pessoal.

\section{API GenJazz}
\label{chap2:sec:api}

A Tabela~\ref{tab:api-endpoints} lista os endpoints utilizados.

\begin{table}[!htb]
\centering\small
\begin{tabular}{llp{5.5cm}}
\toprule
\textbf{Método} & \textbf{Endpoint} & \textbf{Descrição} \\
\midrule
GET & \texttt{/api/keys}, \texttt{/api/structures}, \texttt{/api/modulations} & Catálogos de opções disponíveis \\
GET & \texttt{/api/generate/\{k\}/\{s\}/\{m\}} & Gera uma progressão \\
GET & \texttt{/api/chords2mp3/\{chords\}} & Converte acordes em MP3 \\
GET/POST/DELETE & \texttt{/api/chords/\{email\}/...} & Biblioteca pessoal (listar, guardar, apagar) \\
\bottomrule
\end{tabular}
\caption{Endpoints da API GenJazz.}
\label{tab:api-endpoints}
\end{table}

A biblioteca usa uma estratégia \textit{local-first}: as progressões são guardadas imediatamente no \texttt{localStorage} e a sincronização com a API acontece em segundo plano. Assim a aplicação funciona mesmo sem ligação à internet, com os dados locais disponíveis de imediato.
